% Source: Evans, "Partial Differential Equations" (2nd ed., AMS Graduate Studies in Mathematics vol. 19, 2010),
% Chapter 1 "Introduction", PDF pages 14-26 of the cited source.
% Transcribed from the cited source and organized according to
% leanblueprint conventions (semantic \label, bracketed titles, \uses dependency edges) below.

\section{Partial Differential Equations}
\label{sec:pde-introduction}

\begin{definition}[$k$th-Order Partial Differential Equation]
\label{def:kth-order-pde}
\lean{EvansLib.IsPDESolutionOn}
\leanok
\dcref{ch1:1.1-def-1}
\group{PDE Basics}


Fix an integer $k \geq 1$ and an open set $U \subseteq \R^n$. Given
$$F : \R^{n^k} \times \R^{n^{k-1}} \times \cdots \times \R^n \times \R \times U \to \R,$$
and an unknown $u : U \to \R$, an equation
\begin{equation}
F(D^k u(x), D^{k-1} u(x), \ldots, Du(x), u(x), x) = 0 \qquad (x \in U)
\tag{1}
\end{equation}
is a $k$th-order partial differential equation (PDE).
\end{definition}

A solution satisfies (1), possibly subject to auxiliary boundary conditions on
a subset $\Gamma$ of $\partial U$.

\begin{definition}[Linear, Semilinear, Quasilinear, and Fully Nonlinear PDE]
\label{def:linear-semilinear-quasilinear-fully-nonlinear}
\lean{EvansLib.IsLinearPDE,EvansLib.IsHomogeneousLinearPDE,EvansLib.IsSemilinearPDE,EvansLib.IsQuasilinearPDE,EvansLib.IsFullyNonlinearPDE}
\uses{def:kth-order-pde}
\leanok
\dcref{ch1:1.1-def-2}
\group{PDE Basics}



For (1):

\begin{enumerate}
\item The equation is \emph{linear} if
$$\sum_{|\alpha| \leq k} a_\alpha(x) D^\alpha u = f(x)$$
for given $a_\alpha$ and $f$; it is \emph{homogeneous} when $f \equiv 0$.

\item It is \emph{semilinear} if
$$\sum_{|\alpha|=k} a_\alpha(x) D^\alpha u + a_0(D^{k-1}u, \ldots, Du, u, x) = 0.$$

\item It is \emph{quasilinear} if
$$\sum_{|\alpha|=k} a_\alpha(D^{k-1}u, \ldots, Du, u, x) D^\alpha u + a_0(D^{k-1}u, \ldots, Du, u, x) = 0.$$

\item It is \emph{fully nonlinear} if its dependence on the highest-order derivatives is nonlinear.
\end{enumerate}
\end{definition}

\begin{definition}[$k$th-Order System of Partial Differential Equations]
\label{def:kth-order-system-pde}
\lean{EvansLib.IsPDESystemSolutionOn}
\uses{def:kth-order-pde}
\leanok
\dcref{ch1:1.1-def-3}
\group{PDE Basics}



For an unknown $\mathbf{u}:U\to\R^m$, $\mathbf{u}=(u^1,\ldots,u^m)$, and a given map
$$\mathbf{F} : \R^{mn^k} \times \R^{mn^{k-1}} \times \cdots \times \R^{mn} \times \R^m \times U \to \R^m,$$
a $k^{\text{th}}$-order system is
\begin{equation}
\mathbf{F}(D^k \mathbf{u}(x), D^{k-1}\mathbf{u}(x), \ldots, D\mathbf{u}(x), \mathbf{u}(x), x) = 0 \qquad (x \in U).
\tag{2}
\end{equation}
The displayed system has $m$ scalar equations for the $m$ components of $\mathbf{u}$; systems with fewer or more equations are also allowed.
\end{definition}

\begin{definition}[Classification of Systems]
\label{def:system-linear-semilinear-quasilinear-fully-nonlinear}
\lean{EvansLib.IsLinearPDESystem,EvansLib.IsHomogeneousLinearPDESystem,EvansLib.IsSemilinearPDESystem,EvansLib.IsQuasilinearPDESystem,EvansLib.IsFullyNonlinearPDESystem}
\uses{def:kth-order-system-pde, def:linear-semilinear-quasilinear-fully-nonlinear}
\leanok
\dcref{ch1:1.1-def-2-systems}
\group{PDE Basics}



A system $\mathbf{F}(D^k\mathbf{u}, \ldots, \mathbf{u}, x)=0$ with values in $\R^m$ is classified as in the scalar case. Its coefficients are $\R^m$-valued: it is \emph{linear} if $\mathbf{F} = \sum_{|\alpha|\le k} A_\alpha(x) D^\alpha\mathbf{u} - \mathbf{f}(x)$, with matrix coefficients $A_\alpha$, \emph{semilinear} if the top-order coefficients depend only on $x$, \emph{quasilinear} if they may also depend on lower-order derivatives, and \emph{fully nonlinear} otherwise. Every homogeneous linear system is linear, every linear system is semilinear, every semilinear system is quasilinear, and no quasilinear system is fully nonlinear.
\end{definition}

\begin{definition}[Non-Square and Overdetermined Linear Systems]
\label{def:non-square-linear-system}
\lean{EvansLib.IsLinearPDESystemGen,EvansLib.IsHomogeneousLinearPDESystemGen}
\uses{def:kth-order-system-pde, def:system-linear-semilinear-quasilinear-fully-nonlinear}
\leanok
\group{PDE Basics}



For $\mathbf{u}:U\to\R^m$ and equations valued in $\R^{m'}$, a $k$th-order system is \emph{linear} when
$$\mathbf{F}(D^k\mathbf{u}, \ldots, \mathbf{u}, x) = \sum_{|\alpha|\le k} A_\alpha(x) D^\alpha\mathbf{u} - \mathbf{f}(x),$$
where $A_\alpha(x):\R^m\to\R^{m'}$ are rectangular coefficient maps and $\mathbf{f}:U\to\R^{m'}$; it is \emph{homogeneous} when $\mathbf{f}\equiv0$. The case $m'=m$ recovers \cref{def:system-linear-semilinear-quasilinear-fully-nonlinear}; $m'>m$ is the \emph{overdetermined} case.
\end{definition}

\begin{remark}[``PDE'' as Singular and Plural]
\label{rem:pde-abbreviation}
\dcref{ch1:1.1-rem-1}
\group{PDE Basics}
``PDE'' abbreviates both ``partial differential equation'' and ``partial differential equations''. $\square$
\end{remark}
\section{Examples}
\label{sec:pde-examples}

For $x \in U$, with $U$ open in $\R^n$, and $t \geq 0$, write
$$Du = D_xu = (u_{x_1},\ldots,u_{x_n})$$
for the spatial gradient.

\subsection{Single Partial Differential Equations}
\label{sec:single-pdes}

\subsubsection{Linear Equations}
\label{sec:examples-linear-equations}

\begin{example}[Laplace's Equation]
\label{ex:laplaces-equation}
\lean{EvansLib.laplaceSymbol}
\uses{def:kth-order-pde, def:linear-semilinear-quasilinear-fully-nonlinear}
\leanok
\dcref{ch1:1.2.1.a.1}
\group{Examples of PDE}



$$\Delta u = \sum_{i=1}^n u_{x_i x_i} = 0.$$
\end{example}

\begin{example}[Helmholtz's (or Eigenvalue) Equation]
\label{ex:helmholtz-eigenvalue-equation}
\lean{EvansLib.helmholtzSymbol}
\uses{def:kth-order-pde, def:linear-semilinear-quasilinear-fully-nonlinear}
\leanok
\dcref{ch1:1.2.1.a.2}
\group{Examples of PDE}



$$-\Delta u = \lambda u.$$
\end{example}

\begin{example}[Linear Transport Equation]
\label{ex:linear-transport-equation}
\lean{EvansLib.transportSymbol}
\uses{def:kth-order-pde, def:linear-semilinear-quasilinear-fully-nonlinear}
\leanok
\dcref{ch1:1.2.1.a.3}
\group{Examples of PDE}



$$u_t + \sum_{i=1}^n b^i u_{x_i} = 0.$$
\end{example}

\begin{example}[Liouville's Equation]
\label{ex:liouvilles-equation}
\lean{EvansLib.liouvilleSymbol}
\uses{def:kth-order-pde, def:linear-semilinear-quasilinear-fully-nonlinear}
\leanok
\dcref{ch1:1.2.1.a.4}
\group{Examples of PDE}



$$u_t - \sum_{i=1}^n (b^i u)_{x_i} = 0.$$
\end{example}

\begin{example}[Heat (or Diffusion) Equation]
\label{ex:heat-diffusion-equation}
\lean{EvansLib.heatSymbol}
\uses{def:kth-order-pde, def:linear-semilinear-quasilinear-fully-nonlinear}
\leanok
\dcref{ch1:1.2.1.a.5}
\group{Examples of PDE}



$$u_t - \Delta u = 0.$$
\end{example}

\begin{example}[Schr\"odinger's Equation]
\label{ex:schrodingers-equation}
\lean{EvansLib.schrodingerSystemSymbol}
\uses{def:kth-order-pde, def:system-linear-semilinear-quasilinear-fully-nonlinear}
\leanok
\dcref{ch1:1.2.1.a.6}
\group{Examples of PDE}



$$iu_t + \Delta u = 0.$$
For $u=u^1+iu^2$, the equation is equivalent to the real system
$$\Delta u^1-u^2_t=0,\qquad \Delta u^2+u^1_t=0,$$
or $\Delta\mathbf{u}+i\,\mathbf{u}_t=0$, where $i$ acts by the rotation
$(a,b)\mapsto(-b,a)$. This is a second-order homogeneous linear system.
\end{example}

\begin{example}[Kolmogorov's Equation]
\label{ex:kolmogorovs-equation}
\lean{EvansLib.kolmogorovSymbol}
\uses{def:kth-order-pde, def:linear-semilinear-quasilinear-fully-nonlinear}
\leanok
\dcref{ch1:1.2.1.a.7}
\group{Examples of PDE}



$$u_t - \sum_{i,j=1}^n a^{ij} u_{x_i x_j} + \sum_{i=1}^n b^i u_{x_i} = 0.$$
\end{example}

\begin{example}[Fokker--Planck Equation]
\label{ex:fokker-planck-equation}
\lean{EvansLib.fokkerPlanckSymbol}
\uses{def:kth-order-pde, def:linear-semilinear-quasilinear-fully-nonlinear}
\leanok
\dcref{ch1:1.2.1.a.8}
\group{Examples of PDE}



$$u_t - \sum_{i,j=1}^n (a^{ij} u)_{x_i x_j} - \sum_{i=1}^n (b^i u)_{x_i} = 0.$$
\end{example}

\begin{example}[Wave Equation]
\label{ex:wave-equation-example}
\lean{EvansLib.waveSymbol}
\uses{def:kth-order-pde, def:linear-semilinear-quasilinear-fully-nonlinear}
\leanok
\dcref{ch1:1.2.1.a.9}
\group{Examples of PDE}



$$u_{tt} - \Delta u = 0.$$
\end{example}

\begin{example}[Telegraph Equation]
\label{ex:telegraph-equation}
\lean{EvansLib.telegraphSymbol}
\uses{def:kth-order-pde, def:linear-semilinear-quasilinear-fully-nonlinear}
\leanok
\dcref{ch1:1.2.1.a.10}
\group{Examples of PDE}



$$u_{tt} + d u_t - u_{xx} = 0.$$
\end{example}

\begin{example}[General Wave Equation]
\label{ex:general-wave-equation}
\lean{EvansLib.generalWaveSymbol}
\uses{def:kth-order-pde, def:linear-semilinear-quasilinear-fully-nonlinear}
\leanok
\dcref{ch1:1.2.1.a.11}
\group{Examples of PDE}



$$u_{tt} - \sum_{i,j=1}^n a^{ij} u_{x_i x_j} + \sum_{i=1}^n b^i u_{x_i} = 0.$$
\end{example}

\begin{example}[Airy's Equation]
\label{ex:airys-equation}
\lean{EvansLib.airySymbol}
\uses{def:kth-order-pde, def:linear-semilinear-quasilinear-fully-nonlinear}
\leanok
\dcref{ch1:1.2.1.a.12}
\group{Examples of PDE}



$$u_t + u_{xxx} = 0.$$
\end{example}

\begin{example}[Beam Equation]
\label{ex:beam-equation}
\lean{EvansLib.beamSymbol}
\uses{def:kth-order-pde, def:linear-semilinear-quasilinear-fully-nonlinear}
\leanok
\dcref{ch1:1.2.1.a.13}
\group{Examples of PDE}



$$u_t + u_{xxxx} = 0.$$
\end{example}

\subsubsection{Nonlinear Equations}
\label{sec:examples-nonlinear-equations}

\begin{example}[Eikonal Equation]
\label{ex:eikonal-equation}
\lean{EvansLib.eikonalSymbol}
\uses{def:kth-order-pde, def:linear-semilinear-quasilinear-fully-nonlinear}
\leanok
\dcref{ch1:1.2.1.b.1}
\group{Examples of PDE}



$$|Du| = 1.$$
\end{example}

\begin{example}[Nonlinear Poisson Equation]
\label{ex:nonlinear-poisson-equation}
\lean{EvansLib.nonlinearPoissonSymbol}
\uses{def:kth-order-pde, def:linear-semilinear-quasilinear-fully-nonlinear}
\leanok
\dcref{ch1:1.2.1.b.2}
\group{Examples of PDE}



$$-\Delta u = f(u).$$
\end{example}

\begin{example}[$p$-Laplacian Equation]
\label{ex:p-laplacian-equation}
\lean{EvansLib.pLaplacianSymbol}
\uses{def:kth-order-pde, def:linear-semilinear-quasilinear-fully-nonlinear}
\leanok
\dcref{ch1:1.2.1.b.3}
\group{Examples of PDE}



$$\Div(|Du|^{p-2}Du) = 0.$$
For a classical solution, the equation has the equivalent form
$$|Du|^{p-2}\Delta u + (p-2)|Du|^{p-4}D^2u(Du,Du) = 0.$$
This is a second-order quasilinear equation, and for $p=2$ it becomes Laplace's equation.
\end{example}

\begin{example}[Minimal Surface Equation]
\label{ex:minimal-surface-equation}
\lean{EvansLib.minimalSurfaceSymbol}
\uses{def:kth-order-pde, def:linear-semilinear-quasilinear-fully-nonlinear}
\leanok
\dcref{ch1:1.2.1.b.4}
\group{Examples of PDE}



$$\Div\left(\frac{Du}{(1+|Du|^2)^{1/2}}\right) = 0.$$
For a classical solution, an equivalent form is
$$(1+|Du|^2)^{-1/2}\Delta u - (1+|Du|^2)^{-3/2}D^2u(Du,Du) = 0,$$
and the equation is second-order quasilinear.
\end{example}

\begin{example}[Monge--Amp\`ere Equation]
\label{ex:monge-ampere-equation}
\lean{EvansLib.mongeAmpereSymbol}
\uses{def:kth-order-pde, def:linear-semilinear-quasilinear-fully-nonlinear}
\leanok
\dcref{ch1:1.2.1.b.5}
\group{Examples of PDE}



$$\det(D^2u) = f.$$
\end{example}

\begin{example}[Hamilton--Jacobi Equation]
\label{ex:hamilton-jacobi-equation}
\lean{EvansLib.hamiltonJacobiSymbol}
\uses{def:kth-order-pde, def:linear-semilinear-quasilinear-fully-nonlinear}
\leanok
\dcref{ch1:1.2.1.b.6}
\group{Examples of PDE}



$$u_t + H(Du, x) = 0.$$
This is first-order. If $H$ is nonlinear in its first argument, the equation is
fully nonlinear; in particular, $H(p,x)=|p|^2$ gives $u_t+|Du|^2=0$, whose
symbol is quadratic rather than affine in $Du$.
\end{example}

\begin{example}[Scalar Conservation Law]
\label{ex:scalar-conservation-law}
\lean{EvansLib.conservationLawSymbol}
\uses{def:kth-order-pde, def:linear-semilinear-quasilinear-fully-nonlinear}
\leanok
\dcref{ch1:1.2.1.b.7}
\group{Examples of PDE}



$$u_t + \Div\mathbf{F}(u) = 0.$$
For a classical solution, the equivalent nondivergence form is
$$u_t + \sum_i (F^i)'(u)u_{x_i} = 0.$$
This is a first-order quasilinear equation with top-order coefficients $(F^i)'(u)$.
\end{example}

\begin{example}[Inviscid Burgers' Equation]
\label{ex:inviscid-burgers-equation}
\lean{EvansLib.burgersSymbol}
\uses{def:kth-order-pde, def:linear-semilinear-quasilinear-fully-nonlinear, ex:scalar-conservation-law}
\leanok
\dcref{ch1:1.2.1.b.8}
\group{Examples of PDE}



$$u_t + uu_x = 0.$$
\end{example}

\begin{example}[Scalar Reaction-Diffusion Equation]
\label{ex:scalar-reaction-diffusion-equation}
\lean{EvansLib.reactionDiffusionSymbol}
\uses{def:kth-order-pde, def:linear-semilinear-quasilinear-fully-nonlinear}
\leanok
\dcref{ch1:1.2.1.b.9}
\group{Examples of PDE}



$$u_t - \Delta u = f(u).$$
\end{example}

\begin{example}[Porous Medium Equation]
\label{ex:porous-medium-equation}
\lean{EvansLib.nonlinearDiffusionSymbol}
\uses{def:kth-order-pde, def:linear-semilinear-quasilinear-fully-nonlinear}
\leanok
\dcref{ch1:1.2.1.b.10}
\group{Examples of PDE}



$$u_t - \Delta(u^\gamma) = 0.$$
For a classical solution, an equivalent form is
$$u_t - \gamma u^{\gamma-1}\Delta u - \gamma(\gamma-1)u^{\gamma-2}|Du|^2 = 0,$$
which is second-order quasilinear. The corresponding general nonlinear-diffusion form is
$$u_t - \varphi'(u)\Delta u - \varphi''(u)|Du|^2 = 0$$
with $\varphi(u)=u^\gamma$.
\end{example}

\begin{example}[Nonlinear Wave Equations]
\label{ex:nonlinear-wave-equations}
\lean{EvansLib.nonlinearWaveSymbol,EvansLib.nonlinearWaveDivSymbol}
\uses{def:kth-order-pde, def:linear-semilinear-quasilinear-fully-nonlinear}
\leanok
\dcref{ch1:1.2.1.b.11}
\group{Examples of PDE}



$$u_{tt} - \Delta u = f(u),$$
$$u_{tt} - \Div\,\mathbf{a}(Du) = 0.$$
The reaction form is second-order semilinear. For classical solutions, the
divergence form is equivalent to
$$u_{tt} - \sum_{i,j}\frac{\partial a_j}{\partial p_i}(Du)u_{x_i x_j} = 0,$$
and is second-order quasilinear.
\end{example}

\begin{example}[Korteweg--de Vries (KdV) Equation]
\label{ex:kdv-equation}
\lean{EvansLib.kdvSymbol}
\uses{def:kth-order-pde, def:linear-semilinear-quasilinear-fully-nonlinear}
\leanok
\dcref{ch1:1.2.1.b.12}
\group{Examples of PDE}



$$u_t + uu_x + u_{xxx} = 0.$$
\end{example}

\subsection{Systems of Partial Differential Equations}
\label{sec:systems-of-pdes}

\subsubsection{Linear Systems}
\label{sec:examples-linear-systems}

\begin{example}[Equilibrium Equations of Linear Elasticity]
\label{ex:linear-elasticity-equilibrium}
\lean{EvansLib.linearElasticityEquilibriumSymbol}
\uses{def:kth-order-system-pde, def:system-linear-semilinear-quasilinear-fully-nonlinear}
\leanok
\dcref{ch1:1.2.2.a.1}
\group{Examples of PDE}



$$\mu\Delta\mathbf{u} + (\lambda + \mu)D(\Div \mathbf{u}) = 0.$$
This is a second-order linear system for $\mathbf{u}:\R^n\to\R^n$.
\end{example}

\begin{example}[Evolution Equations of Linear Elasticity]
\label{ex:linear-elasticity-evolution}
\lean{EvansLib.linearElasticityEvolutionSymbol}
\uses{def:kth-order-system-pde, def:system-linear-semilinear-quasilinear-fully-nonlinear, ex:linear-elasticity-equilibrium}
\leanok
\dcref{ch1:1.2.2.a.2}
\group{Examples of PDE}



$$\mathbf{u}_t - \mu\Delta\mathbf{u} - (\lambda + \mu)D(\Div \mathbf{u}) = 0.$$
This is a second-order linear system.
\end{example}

\begin{example}[Maxwell's Equations]
\label{ex:maxwells-equations}
\lean{EvansLib.maxwellSymbol,EvansLib.maxwellSymbol_isLinearPDESystemGen}
\uses{def:kth-order-system-pde, def:non-square-linear-system}
\leanok
\dcref{ch1:1.2.2.a.3}
\group{Examples of PDE}



$$\begin{cases}
\mathbf{E}_t = \curl \mathbf{B} \\
\mathbf{B}_t = -\curl \mathbf{E} \\
\Div \mathbf{B} = \Div \mathbf{E} = 0.
\end{cases}$$
For $\mathbf{u}=(\mathbf{E},\mathbf{B}):\R^4\to\R^6$, with
$\mathbf{E}=(u^0,u^1,u^2)$ and $\mathbf{B}=(u^3,u^4,u^5)$, this is a
first-order homogeneous linear system with $8$ equations for $6$ unknowns:
the curl equations contribute $3+3$ equations and the divergence constraints
contribute $1+1$. Thus it is overdetermined, with symbol valued in $\R^8$.
\end{example}

\subsubsection{Nonlinear Systems}
\label{sec:examples-nonlinear-systems}

\begin{example}[System of Conservation Laws]
\label{ex:system-conservation-laws}
\lean{EvansLib.conservationLawSystemSymbol}
\uses{def:kth-order-system-pde, def:system-linear-semilinear-quasilinear-fully-nonlinear, ex:scalar-conservation-law}
\leanok
\dcref{ch1:1.2.2.b.1}
\group{Examples of PDE}



$$\mathbf{u}_t + \Div \mathbf{F}(\mathbf{u}) = 0.$$
For a classical solution, the equivalent nondivergence form is
$$\mathbf{u}_t + \sum_i (D\mathbf{F}^i)(\mathbf{u})\,\mathbf{u}_{x_i} = 0.$$
This is a first-order quasilinear system with top-order coefficient matrices
$D\mathbf{F}^i(\mathbf{u})$.
\end{example}

\begin{example}[Reaction-Diffusion System]
\label{ex:reaction-diffusion-system}
\lean{EvansLib.reactionDiffusionSystemSymbol}
\uses{def:kth-order-system-pde, def:system-linear-semilinear-quasilinear-fully-nonlinear, ex:scalar-reaction-diffusion-equation}
\leanok
\dcref{ch1:1.2.2.b.2}
\group{Examples of PDE}



$$\mathbf{u}_t - \Delta\mathbf{u} = \mathbf{f}(\mathbf{u}).$$
This is a second-order semilinear system with lower-order nonlinearity
$\mathbf{f}(\mathbf{u})$.
\end{example}

\begin{example}[Euler's Equations for Incompressible, Inviscid Flow]
\label{ex:euler-equations-incompressible-inviscid}
\lean{EvansLib.eulerSymbol,EvansLib.eulerSymbol_isQuasilinearPDESystem}
\uses{def:kth-order-system-pde, def:system-linear-semilinear-quasilinear-fully-nonlinear}
\leanok
\dcref{ch1:1.2.2.b.3}
\group{Examples of PDE}



$$\begin{cases}
\mathbf{u}_t + \mathbf{u} \cdot D\mathbf{u} = -D p \\
\Div \mathbf{u} = 0.
\end{cases}$$
The pair $(\mathbf{u},p):\R^{n+1}\to\R^{n+1}$ satisfies the $(n+1)\times(n+1)$ first-order system
$$u^i_t+\sum_j u^j u^i_{x_j}+p_{x_i}=0,\qquad \sum_j u^j_{x_j}=0.$$
This is a first-order quasilinear system with coefficient $u^j$ on
$u^i_{x_j}$.
\end{example}

\begin{example}[Navier--Stokes Equations for Incompressible, Viscous Flow]
\label{ex:navier-stokes-equations-incompressible-viscous}
\lean{EvansLib.nsSymbol,EvansLib.nsSymbol_isSemilinearPDESystem}
\uses{def:kth-order-system-pde, def:system-linear-semilinear-quasilinear-fully-nonlinear, ex:euler-equations-incompressible-inviscid}
\leanok
\dcref{ch1:1.2.2.b.4}
\group{Examples of PDE}



$$\begin{cases}
\mathbf{u}_t + \mathbf{u} \cdot D\mathbf{u} - \Delta\mathbf{u} = -D p \\
\Div \mathbf{u} = 0.
\end{cases}$$
This is obtained from Euler's equations by adding $-\Delta\mathbf{u}$. It is a
second-order semilinear system whose lower-order part is Euler's operator.
\end{example}

See \cite{Zwillinger1984} for further examples of partial differential equations.
\section{Strategies for Studying PDE}
\label{sec:strategies-studying-pde}

\subsection{Well-Posed Problems, Classical Solutions}
\label{sec:well-posed-classical-solutions}

\begin{definition}[Well-Posed Problem]
\label{def:well-posed-problem}
\lean{EvansLib.PDEProblem.IsWellPosed}
\leanok
\dcref{ch1:1.3.1-def-1}
\group{Well-Posedness}


A problem for a partial differential equation is \textit{well-posed} when it has a solution, that solution is unique, and the solution depends continuously on the prescribed data.
\end{definition}

\begin{definition}[Classical Solution]
\label{def:classical-solution}
\lean{EvansLib.IsClassicalSolutionOn}
\uses{def:kth-order-pde}
\leanok
\dcref{ch1:1.3.1-def-2}
\group{Well-Posedness}



For a PDE of order $k$, a \textit{classical solution} is a solution that is at least $k$ times continuously differentiable, so that all derivatives occurring in the equation exist and are continuous.
\end{definition}

\subsection{Weak Solutions and Regularity}
\label{sec:weak-solutions-regularity}

For the scalar conservation law
$$u_t + F(u)_x = 0$$
classical solutions may fail after shocks form, while a suitable class of
\emph{weak} or \emph{generalized} solutions can restore well-posedness.

Existence, uniqueness, and continuous dependence in a weak solution class are
separate from regularity. A regular weak solution is classical; otherwise the
weak solution remains the relevant solution concept.

\section{Multiindex Calculus}
\label{sec:multiindex-calculus}

%ORIGINAL-NUMBER: 2
\begin{proposition}[Multinomial Theorem]
\label{ex:multinomial-theorem}
\lean{EvansLib.multinomial_theorem}
\leanok
\group{Multiindex Calculus}


For $x = (x_1, \ldots, x_n)$ and $k \in \N$,
\[
\left(x_1 + \cdots + x_n\right)^k
= \sum_{|\alpha|=k} \binom{|\alpha|}{\alpha}x^\alpha,
\]
where $\binom{|\alpha|}{\alpha} := |\alpha|!/\alpha!$,
$\alpha! := \alpha_1!\cdots\alpha_n!$, $x^\alpha :=
x_1^{\alpha_1}\cdots x_n^{\alpha_n}$, and the sum ranges over
$\alpha \in \N^n$ with $|\alpha|=k$.
\end{proposition}
\begin{proof}
\leanok
\sketch
Expand the $k$ factors and index each resulting monomial by a function
$f : \{1,\ldots,k\} \to \{1,\ldots,n\}$. Setting
$\alpha_i := \#\{r : f(r)=i\}$ gives $|\alpha|=k$ and the monomial
$x^\alpha$. For fixed $\alpha$, exactly
$k!/ (\alpha_1!\cdots\alpha_n!) = \binom{|\alpha|}{\alpha}$ functions
produce it. Grouping by $\alpha$ proves the identity.
\end{proof}

\begin{definition}[Partial Derivative]
\label{def:partial-derivative}
\lean{EvansLib.partialDeriv}
\leanok
\group{Multiindex Calculus}


For $f : \R^n \to \R$ and $i \in \{1,\ldots,n\}$,
\[
\partial_{x_i}f(x) := Df(x)(e_i),
\]
where $e_i$ is the $i$th standard basis vector.
\end{definition}

\begin{definition}[Multiindex Partial Derivative]
\label{def:multiindex-partial-derivative}
\lean{EvansLib.multiPartial}
\uses{def:partial-derivative}
\leanok
\group{Multiindex Calculus}



For $\alpha=(\alpha_1,\ldots,\alpha_n)\in\N^n$ and
$f : \R^n \to \R$,
\[
D^\alpha f := \frac{\partial^{|\alpha|}f}
 {\partial x_1^{\alpha_1}\cdots\partial x_n^{\alpha_n}}
= \partial_{x_1}^{\alpha_1}\cdots\partial_{x_n}^{\alpha_n}f,
\qquad |\alpha|=\alpha_1+\cdots+\alpha_n.
\]
\end{definition}

\begin{proposition}[Equality of Mixed Partial Derivatives]
\label{prop:clairaut-partial-commute}
\lean{EvansLib.partialDeriv_comm}
\uses{def:partial-derivative, def:multiindex-partial-derivative}
\leanok
\group{Multiindex Calculus}



For smooth $f : \R^n \to \R$ and indices $i,j$,
\[
\partial_{x_i}\partial_{x_j}f=\partial_{x_j}\partial_{x_i}f.
\]
Thus $D^\alpha f$ is independent of the order of the coordinate partials.
\end{proposition}
\begin{proof}
\leanok
\sketch
For a $C^2$ function,
\[
\partial_{x_i}\partial_{x_j}f(x)=D^2f(x)(e_i,e_j).
\]
The Clairaut--Schwarz theorem makes $D^2f(x)$ symmetric, yielding the
commutation identity. Any two orderings differ by adjacent transpositions,
so the multiindex derivative is well defined.
\end{proof}

\begin{proposition}[Iterated Leibniz Rule Along One Axis]
\label{prop:iterated-leibniz-one-axis}
\lean{EvansLib.partialDeriv_iterate_mul}
\uses{def:partial-derivative}
\leanok
\group{Multiindex Calculus}



For smooth $u,v : \R^n \to \R$, an axis $i$, and $m\in\N$,
\[
\partial_{x_i}^m(uv)=\sum_{j=0}^m\binom{m}{j}
(\partial_{x_i}^j u)(\partial_{x_i}^{m-j}v).
\]
\end{proposition}
\begin{proof}
\leanok
\sketch
Restrict to the line $\gamma(t)=x+te_i$. The chain rule identifies the
$k$th derivatives of the restrictions of $u$, $v$, and $uv$ with the
$k$th $x_i$-derivatives. Applying the one-variable Leibniz rule at $t=0$
gives the formula at $x$.
\end{proof}

%ORIGINAL-NUMBER: 3
\begin{proposition}[Leibniz' Formula]
\label{prop:leibniz-formula}
\lean{EvansLib.multiPartial_mul}
\uses{def:multiindex-partial-derivative, prop:clairaut-partial-commute, prop:iterated-leibniz-one-axis}
\leanok
\group{Multiindex Calculus}



For smooth $u,v : \R^n \to \R$,
\[
D^\alpha(uv)=\sum_{\beta\leq\alpha}\binom{\alpha}{\beta}
D^\beta u\,D^{\alpha-\beta}v,
\]
where
$\binom{\alpha}{\beta}:=\alpha!/\bigl(\beta!(\alpha-\beta)!\bigr)
=\prod_{i=1}^n\binom{\alpha_i}{\beta_i}$ and
$\beta\leq\alpha$ means $\beta_i\leq\alpha_i$ for every $i$.
\end{proposition}
\begin{proof}
\leanok
\sketch
Induct on $|\alpha|$. The case $|\alpha|=0$ is immediate. If
$\alpha_i\geq1$, write $D^\alpha=\partial_{x_i}D^{\alpha-e_i}$ using
\cref{prop:clairaut-partial-commute}, apply the induction hypothesis, and
differentiate each product. The coefficient of
$D^\beta u\,D^{\alpha-\beta}v$ is
\[
\binom{\alpha-e_i}{\beta-e_i}+\binom{\alpha-e_i}{\beta},
\]
which equals $\binom{\alpha}{\beta}$ by Pascal's rule in the $i$th factor.
\end{proof}

%ORIGINAL-NUMBER: 4
\begin{proposition}[Directional Derivatives via the Multinomial Theorem]
\label{prop:directional-derivative-multinomial}
\lean{EvansLib.iteratedDeriv_scaled_line_multinomial}
\uses{def:multiindex-partial-derivative, prop:clairaut-partial-commute, ex:multinomial-theorem}
\leanok
\group{Multiindex Calculus}



For smooth $f : \R^n \to \R$, fixed $x\in\R^n$, and
$g(t):=f(tx)$, every $j\in\N$ satisfies
\[
g^{(j)}(0)=\sum_{|\alpha|=j}\binom{j}{\alpha}D^\alpha f(0)x^\alpha,
\]
where $\binom{j}{\alpha}:=j!/\alpha!$ and
$x^\alpha=x_1^{\alpha_1}\cdots x_n^{\alpha_n}$.
\end{proposition}
\begin{proof}
\leanok
\sketch
Define $D_xh(y):=Dh(y)(x)=\sum_i x_i\partial_{x_i}h(y)$. The chain rule
gives $g^{(j)}(t)=(D_x^j f)(tx)$. Since coordinate partials commute by
\cref{prop:clairaut-partial-commute}, the Multinomial Theorem yields
\[
D_x^j=\left(\sum_i x_i\partial_{x_i}\right)^j
=\sum_{|\alpha|=j}\binom{j}{\alpha}x^\alpha D^\alpha.
\]
Evaluating at $0$ proves the result; the operator identity follows by
induction, with Pascal's relation recombining the terms indexed by
$\alpha+e_i$.
\end{proof}

\begin{proposition}[Taylor's Formula in Multiindex Notation]
\label{prop:taylor-multiindex}
\lean{EvansLib.taylor_multiindex}
\uses{def:multiindex-partial-derivative, prop:directional-derivative-multinomial}
\leanok
\group{Multiindex Calculus}



For smooth $f : \R^n \to \R$ and $k=1,2,\ldots$,
\[
f(x)=\sum_{|\alpha|\leq k}\frac{1}{\alpha!}D^\alpha f(0)x^\alpha
+O(|x|^{k+1})\qquad\text{as }x\to0.
\]
\end{proposition}
\begin{proof}
\leanok
\sketch
Set $g(t)=f(tx)$. Taylor's theorem at $0$ with Lagrange remainder gives
\[
g(1)=\sum_{j=0}^k\frac{g^{(j)}(0)}{j!}
 +\frac{g^{(k+1)}(\xi)}{(k+1)!},\qquad \xi\in(0,1).
\]
By \cref{prop:directional-derivative-multinomial},
$g^{(j)}(0)/j!=\sum_{|\alpha|=j}D^\alpha f(0)x^\alpha/\alpha!$.
For the remainder, the same formula at $\xi x$ expresses
$g^{(k+1)}(\xi)$ as a sum of terms
$\binom{k+1}{\alpha}D^\alpha f(\xi x)x^\alpha$ with $|\alpha|=k+1$.
The derivatives are bounded by some $M$ on the closed unit ball, and for
$|x|\leq1$ one has $\xi x$ in that ball and
$|x^\alpha|\leq|x|^{k+1}$. Hence the remainder is
$O(|x|^{k+1})$.
\end{proof}
